\Chapter{49}

\Verse{1}
{
	\hJ{וַיִּקְרָא יַעֲקֹב אֶל בָּנָיו וַיֹּאמֶר}
	\hJ{הֵאָסְפוּ וְאַגִּידָה לָכֶם אֵת אֲשֶׁר יִקְרָא אֶתְכֶם בְּאַחֲרִית הַיָּמִים׃}
}
{
	\eJ{
		And Jacob called to his sons and said:
	}
	\eJ{
		Gather, and I’ll tell you what will happen to you in the future days.
	}
}
{
	Strap in folks.

	Here we have the Blessing of Jacob.

	It's one of the oldest things in the bible.

	There's no way to explain what's going on here.

	... We can't just say some words and color the sources.

	... The book hasn't had to do something like this before.

	... Not sure how to make this work, but w/e, fuck it, follow me.


	\aB{
		This needs one bit of ground rules before the surgery starts. ``MT'' means the Masoretic Text:
		the old Hebrew consonantal text as preserved in the medieval Masoretic tradition, plus the much
		later vowel points and reading tradition. The consonants are ancient, but this poem is older
		still, and some of its words were already weird by the time our surviving witnesses were copied.
		We do have bits of Gen. 49 at Qumran: 4QGen-Exod\textsuperscript{a} preserves parts of vv. 1--5
		and 4QGen\textsuperscript{e} parts of vv. 6--8. They are broadly MT-like; nobody found a Dead
		Sea Scroll containing the neat reconstructed poem printed below. The reconstruction in this file
		is much more aggressive than a normal Bible translation: it uses old spelling, meter,
		repeated-letter mistakes, parallel poems, the Samaritan Pentateuch, and ancient translations to
		guess backward. Angle brackets mean something has been supplied; empty brackets mark something
		being removed; long blanks mean, basically, ``we don't know.'' The LXX is especially useful
		because its Greek was translated from Hebrew centuries before the medieval MT manuscripts---but
		Greek can preserve an older reading, or just an ancient translator's best guess. So below I'll
		keep saying which kind of evidence we've actually got.
	}
}


\Verse{1a}
{
	\hJ{וַיִּקְרָא יַעֲקֹב אֶל בָּנָיו וַיֹּאמֶר}
}
{
	\eJ{
		And Jacob called to his sons and said:
	}
}
{
    This is the only part that's not a mess.

    May as well use this opportunity to pause and take a breath.

    \aC{
        We take up a distinctly different genre of ancient Yahwistic poetry with the
        tribal blessings, the Blessing of Jacob and the Blessing of Moses. To the same
        genre belong the tribal references in Jud. 5:14--18 and Ps. 68:28.\textsuperscript{a}
        The Blessings also have certain characteristics in common with the Oracles of
        Balaam, but there are also strong contrasts.\textsuperscript{b}

        It seems probable in the light of \emph{gattungs- und überlieferungs-geschichtlich}
        studies that individual blessings rest on ancient tradition. They circulated in
        oral form as folk literature, finally being gathered into collections. Some
        blessings are composite, such as those of Judah (Gen. 49:8--12) and Levi
        (Deut. 33:8--11), made up of originally separate elements. These illustrate the
        way in which blessings were brought together. On the other hand, the tendency
        to chop up the blessings into a host of fragments is the result of overly
        atomistic analysis. The Blessing of Jacob consists of a nucleus of blessings
        consciously set in the mouth of the Patriarch (e.g., Reuben, Judah), to which
        have been appended blessings drawn from other sources (e.g., the Joseph
        blessing, which occurs in slightly altered form in the Blessing of Moses). We
        may assume that groups of blessings, ascribed to Jacob and Moses, and perhaps
        others, circulated orally in the period of the Judges.

        After the earlier classical pieces, the Songs of Deborah and Miriam, the
        blessings in these two collections are disappointing as literature. They are
        practically folk poetry, and lack the varied meter and complex rhetorical style
        of the victory odes. They delight in plays on words, frequently use rhyme, and
        with certain notable exceptions are uninspired in content.

        With our increasing knowledge of the amphictyonic organization of the tribes
        in the period of the Judges,\textsuperscript{c} the place of the collection of
        tribal blessings in the religious and political life of the people has become
        clear. One may be assured that such collections were framed in this period from
        various cycles of oral tradition. It is most unlikely that the blessings were
        artificially set in a framework after the ``twelve''-tribe system had ceased to
        be a living organization.

        No serious question has been raised by scholars as to the pre-monarchic date of
        the majority of the individual blessings in Gen. 49. The only exception worthy
        of comment is the blessing of Judah, Gen. 49:8--11. As will be pointed out in
        the study which follows, this blessing is composite: vss. 8, 9, and 11--12 are
        individual units of a most archaic type. The short enigmatic blessing, vs. 10,
        is usually referred to the Judahite monarchy, the earliest possibility being
        the reign of David. An allusion to David is possible; more probably we have
        here pre-Davidic material which came naturally to be associated with the great
        king, was interpreted and may have been slightly modified in the light of his
        career and achievements. This is speculative, however, and it is quite
        impossible to date the collection on the basis of a doubtful understanding of
        a corrupt text. Judahite hegemony in the south during the period of the Judges
        may be all that is implied. It may be remarked that certain blessings (e.g.,
        Reuben, Simeon and Levi) may go back to patriarchal traditions. It is best to
        date the collection in substantially its present form in the late period of the
        Judges. Orthographic and linguistic criteria support such a date.
    }


	\aB{
		This little prose frame is the easy part. MT has ויקרא יעקב אל בניו ויאמר (\emph{wayyiqra
		Yaʿaqov el banaw wayyomer}), ``Jacob called to his sons and said,'' and there is no serious
		textual reason to repair it. The early Qumran copy 4QGen-Exod\textsuperscript{a} reaches the
		opening of this chapter and gives us no reason to think some radically different opening once
		stood here. The real mess starts when the poem itself begins.
	}
}

\Verse{1b}
{
	האסף ואגד לכם [ ]\textsuperscript{1}
}
{
	Gather yourselves that I may inform you,
}
{
	\aB{
		MT keeps going after this with את אשר יקרא אתכם באחרית הימים (\emph{et asher yiqra etkhem
		beʾaḥarit hayyamim}), ``what will happen to you in the latter days.'' This reconstruction throws
		that whole tail away as a prose expansion and keeps only ``Gather, and I'll tell you.''
		Important: that longer wording is also in the LXX---τί ἀπαντήσει ὑμῖν ἐπ᾽ ἐσχάτων τῶν ἡμερῶν
		(\emph{ti apantēsei hymin ep' eschatōn tōn hēmerōn}), basically ``what will meet you in the last
		days.'' So this is \emph{not} a lost old line recovered from Greek. It is a literary and
		metrical guess that the shorter poem came first and somebody later made Jacob's introduction
		more explicit.
	}
}

\Verse{2}
{
	הקבצו [ ] בני יעקב

	ושמע אל ישראל אבכם
}
{
	Assemble together, O sons of Jacob,

	Hearken unto Israel your father.
}
{
	\aB{
		MT says הקבצו ושמעו בני יעקב ושמעו אל ישראל אביכם (\emph{hiqqavetsu weshimʿu bene Yaʿaqov;
		weshimʿu el Yisraʾel avikhem}), with ``listen'' twice. The reconstruction deletes the first ושמעו (\emph{weshimʿu}) and makes vv. 1b--2 into three short balanced lines: gather / sons of
		Jacob / listen to Israel your father. The idea is simple scribal copying: a scribe's eye or ear
		picked up the ושמעו from the next line and duplicated it. But the ancient versions also know the
		repeated ``listen,'' so this is mostly a meter-based emendation, not something another
		manuscript hands us on a plate.
	}
}

\Verse{3}
{
	ראובן בכר

	את\textsuperscript{2} כח

	[ו־]\textsuperscript{3} ראשית אן

	[ ]\textsuperscript{4} יתר עז\textsuperscript{5}\textsuperscript{6}
}
{
	Reuben, my first-born

	Thou art my strength

	The prime of my manhood

	The best of my might.
}
{
	\aB{
		MT is crowded here: ראובן בכרי אתה כחי וראשית אוני יתר שאת ויתר עז (\emph{Reʾuven bekhori atta,
		koḥi wereshit oni, yeter seʾet weyeter ʿaz}). The reconstruction turns it back into four tiny
		rhyming lines. It moves אתה (\emph{atta}, ``you'') with ``my strength,'' drops the troublesome יתר שאת (\emph{yeter seʾet}, traditionally something like ``excellence of dignity''), and keeps יתר עז but reads the last word as עזי (\emph{ʿuzzi}, ``my might'') with the final \emph{-i}
		understood rather than written. That missing final vowel-letter fits genuinely old Hebrew
		spelling, and it makes all four cola end in ``my ---.'' The bigger deletion, though, is
		conjectural: the endnote thinks יתר שאת grew out of an accidental repeat or corruption around ראשית (\emph{reshit}, ``first/beginning''). The versions show that ancient readers already
		struggled with the end of this verse; they do not preserve one clean alternate Hebrew text.
	}
}

\Verse{4}
{
	\_\_\_\_\_\_

	\textsuperscript{7}\_\_\_\_\_\_

	כ עלת <על>\textsuperscript{8} משכבי\textsuperscript{9} [ ]

	אז חללת יצעי\textsuperscript{10} <אבכם>\textsuperscript{11}
}
{
	\ldots\ldots\ldots

	\ldots\ldots\ldots

	For thou didst go up to my bed

	Yea, thou didst profane the couch
	of thy father.
}
{
	\aB{
		The first half of MT is פחז כמים אל תותר (\emph{paḥaz kammayim al totar}), usually made into
		something like ``unstable/reckless as water, you won't excel.'' The endnote thinks the line is
		too damaged to recover and the file therefore prints blanks instead of pretending. The LXX
		already has a different attempt at sense: ἐξύβρισας ὡς ὕδωρ, μὴ ἐκζέσῃς (\emph{exybrisas hōs
		hydōr, mē ekzēsēs}), roughly ``you ran wild like water; don't boil over.'' That tells us the
		line was hard very early, but it does not give us the lost Hebrew. The second half is more
		repairable. After עלית (\emph{ʿalita}, ``you went up'') the reconstruction supplies על
		(\emph{ʿal}, ``onto''), reads משכבי (\emph{mishkavi}) as singular ``my bed,'' and moves ``your
		father'' so the two lines become, plainly, ``you went up to my bed / you profaned your father's
		couch.'' The parallel wording in 1 Chr. 5:1 is the main external clue for that word move.
	}
}

\Verse{5}
{
	שמעון ולו אחים\textsuperscript{12}

	כל חמס מכרתם [ ]\textsuperscript{13}
}
{
	Simeon and Levi are brethren

	Weapons of violence are their merchandise.
}
{
	\aB{
		The problem is mostly one horrible word: מכרתיהם (\emph{mekherotehem}), which occurs nowhere
		else in the Bible. MT's consonants are there; we just do not really know what the word means.
		Traditional translations make it some kind of ``weapon,'' while the LXX goes off in another
		direction entirely: συνετέλεσαν ἀδικίαν ἐξ αἱρέσεως αὐτῶν (\emph{synetelesan adikian ex
		haireseōs autōn}), roughly ``they carried out injustice by their own choice.'' This
		reconstruction instead connects the word with the Northwest Semitic root מ־כ־ר (\emph{m-k-r}),
		``trade/sell,'' including Ugaritic evidence, and gets ``violence is their stock in trade /
		merchandise.'' It also prefers a shorter old suffix, \emph{-am}. So nothing big is physically
		missing from MT here; the ancient word has simply become opaque, and this translation is one
		philological bet among several.
	}
}

\Verse{6}
{
	בסדם אל תבא\textsuperscript{14} נפש

	בקהלם אל תחד\textsuperscript{15} כבדי\textsuperscript{16}

	[ ]\textsuperscript{17} באפם הרג אש

	[ו־]\textsuperscript{18} ברצנם עקר שר\textsuperscript{19}
}
{
	Into their council enter not, O my soul

	In their assembly join not, O my heart.

	In their anger, they slew a man

	Willfully they houghed a bull.
}
{
	\aB{
		Here there is one especially nice old-version clue. MT has כבודי (\emph{kevodi}), normally ``my
		glory,'' but this reconstruction reads כבדי (\emph{kevedi}), ``my liver.'' That sounds bizarre
		to us, but the liver was an ancient seat of feeling, so ``my soul'' and ``my liver'' make a
		normal poetic pair. The LXX actually has τὰ ἥπατά μου (\emph{ta hēpata mou}), ``my liver(s),''
		which is strong evidence that its Hebrew Vorlage was read that way; Ugaritic also pairs
		\emph{kbd}, ``liver,'' with other emotion/body words. The file then reads תבא (\emph{tavo}) as
		an old spelling of feminine תבאי (\emph{tavoʾi}), because Jacob is talking directly to ``my
		soul,'' and understands תחד (\emph{teḥad}) as ``join [with them],'' with ``with them'' left
		unstated. It drops כי (\emph{ki}) and a line-opening \emph{waw} for symmetry. Qumran preserves
		only scraps here, including בקהלם (\emph{biqhalam}, ``in their assembly''), so it does not
		settle these little readings.
	}
}

\Verse{7}
{
	ארר אפם כי עז

	ועברתם כי קשת

	אחלקם\textsuperscript{20} ביעקב

	אפצם\textsuperscript{21} בישראל
}
{
	Cursed be their anger, for it is fierce

	And their wrath, for it is harsh.

	I will divide them in Jacob,

	I will scatter them in Israel.
}
{
	\aB{
		This one is basically healthy. The reconstruction is mostly showing what the same poem could
		look like in much older spelling: ארר for later ארור (\emph{arur}, ``cursed'') and אפצם for
		later אפיצם (\emph{afitsem}, ``I will scatter them''), with vowel letters not yet written. The
		first-person ``I will divide them / I will scatter them'' is not treated as corruption; Jacob is
		speaking the tribal fate as a curse. 4QGen\textsuperscript{e} preserves part of this verse and
		is broadly compatible with MT. So this is a good reminder that weird-looking spelling in the
		reconstructed text does not always mean the MT lost words.
	}
}

\Verse{8}
{
	יהד הד\textsuperscript{22} את\textsuperscript{23}

	ידך\textsuperscript{24} בערף איבך

	יהד ידך אחך\textsuperscript{25}

	ישתחו לך בני אבך\textsuperscript{25a}
}
{
	Judah, my splendor art thou,

	Thy hand is on the neck of thine enemies.

	Judah, thy brothers shall praise thee

	The sons of thy father shall bow
	down to thee.
}
{
	\aB{
		This reconstruction thinks the beginning got scrambled because several very short old spellings
		looked almost alike. In early spelling יהודה (\emph{Yehuda}) could be יהד, while a restored הודי
		(\emph{hodi}, ``my splendor/majesty'') could be הד. The file therefore reconstructs ``Judah, my
		splendor are you,'' moves ידך בערף איבך (\emph{yadkha beʿoref oyevekha}, ``your hand on your
		enemies' neck'') next to it, and then gives ``Judah, your brothers praise you.'' The ancient
		versions help with the last line: LXX has Ἰούδα, σὲ αἰνέσαισαν οἱ ἀδελφοί σου (\emph{Iouda, se
		ainesaisan hoi adelphoi sou}), ``Judah, may your brothers praise you,'' without anything that
		requires MT's separate nominative אתה (\emph{atta}, ``you'') in that clause. But the restored הודי itself is still a guess based on meter and likely eye-skip. Qumran preserves the later
		words בערף אויביך ... ישתחוו, ``on the neck of your enemies ... they will bow,'' but not enough
		to prove the new first line. So: clever reconstruction, not recovered manuscript text.
	}
}

\Verse{9}
{
	גר ארי יהד\textsuperscript{26}

	מטרף בני עלת\textsuperscript{27}

	כרע רבץ כארי\textsuperscript{28}

	[ו־]\textsuperscript{29} כלבא מי יקמנו\textsuperscript{30}
}
{
	A lion's whelp is Judah

	From the prey, my son, thou hast gone up.

	He crouches, he couches like a lion,

	Like a lioness, who dares rouse him?
}
{
	\aB{
		There is very little textual surgery here. MT already gives the lion poem: ``Judah is a lion cub
		... he crouched, lay down like a lion ... who will raise him?'' The file mostly strips later
		vowel letters and drops the \emph{waw} before the final colon because very old Canaanite/Hebrew
		poetry often begins cola without ``and.'' The phrase מטרף בני עלית (\emph{mitteref beni
		ʿalita}), ``from the prey, my son, you went up,'' is still a little cryptic, but that is an
		interpretation problem, not evidence that a line is missing. The LXX has the same basic lion
		picture, so there is no useful alternate text to import.
	}
}

\Verse{10}
{
	לא יסר שפט\textsuperscript{31} מיהד

	ומחקק\textsuperscript{32} מבין\textsuperscript{33} דגלו

	\_\_\_\_\_\_

	\textsuperscript{34}\_\_\_\_\_\_
}
{
	There shall not fail a judge from Judah

	Nor a commander from
	among his standards

	\ldots\ldots\ldots

	\ldots\ldots\ldots
}
{
	\aB{
		This is one of the giant ones. MT begins לא יסור שבט מיהודה ומחקק מבין רגליו (\emph{lo yasur
		shevet mi-Yehuda umeḥoqeq mibben raglaw}), usually ``the scepter won't depart from Judah, nor
		the ruler's staff from between his feet.'' The file changes שבט (\emph{shevet},
		``staff/scepter'') to שפט (\emph{shofet/shefet}, ``judge, ruler''): old \emph{bet} and \emph{pe}
		were easy to confuse, and LXX has ἄρχων (\emph{archōn}), ``ruler.'' מחקק (\emph{meḥoqeq}) is
		then ``commander.'' Next it follows the Samaritan Pentateuch's דגליו (\emph{degalaw}, ``his
		standards/banners'') instead of MT רגליו (\emph{raglaw}, ``his feet/legs''). That gives ``a
		commander from among his banners.'' And yes: MT's ``between his legs'' absolutely can be a
		genital/crotch euphemism; this reconstruction is specifically choosing a textual variant that
		makes the dick-reading disappear. The second half is worse. MT's עד כי יבא שילה ולו יקהת עמים
		(\emph{ʿad ki yavo Shilo welo yiqqehat ʿammim}) is famously opaque. LXX has ἕως ἂν ἔλθῃ τὰ
		ἀποκείμενα αὐτῷ ... προσδοκία ἐθνῶν (\emph{heōs an elthē ta apokeimena autō ... prosdokia
		ethnōn}), something like ``until what is laid up for him comes ... [he is] the expectation of
		nations.'' The endnote can suggest a tribute/submission reconstruction, but it cannot prove it,
		so this file does the sensible thing and leaves the whole bicolon blank.
	}
}

\Verse{11}
{
	אסר\textsuperscript{35} לגפנו\textsuperscript{36} ערו\textsuperscript{37}

	[ו־]\textsuperscript{38} לשרק\textsuperscript{39} בני\textsuperscript{40} אתנו

	כבס בין לבש

	[ו־]\textsuperscript{41} בדם ענבים סות\textsuperscript{42}
}
{
	His ass is harnessed with its bridle (?)

	The colt of his she-ass with its trappings (?)

	He dyes with wine his garment

	With the blood of grapes his mantle.
}
{
	\aB{
		The first two lines are a different kind of reconstruction: mostly the same consonants, but read
		with a lost old vocabulary. MT vocalizes גפן (\emph{gefen}) as ``vine'' and שרקה (\emph{soreqa})
		as ``choice vine,'' giving the famous absurdly-rich vineyard picture: tie the donkey to a vine
		because grapes are basically weeds now. This file instead uses Ugaritic parallels to read the
		words as horse/donkey gear---something like ``his ass is harnessed with its bridle / the colt
		... with its trappings.'' That is not supplied by another Hebrew manuscript. In fact LXX plainly
		follows the vineyard reading: δεσμεύων πρὸς ἄμπελον τὸν πῶλον αὐτοῦ (\emph{desmeuōn pros ampelon
		ton pōlon autou}), ``binding his colt to a vine.'' So this first bicolon is a bold lexical
		re-reading, not a secure textual correction. The wine-and-clothes half is much firmer. The rare סות (\emph{sut}), ``mantle/garment,'' has good Phoenician parallels, and the file mainly removes
		line-opening \emph{waw}s and reads the old suffix spellings consistently.
	}
}

\Verse{12}
{
	חכלל\textsuperscript{43} ענם מיין

	לבן\textsuperscript{44} שנם מחלב\textsuperscript{45}
}
{
	Darker are his eyes than wine

	Whiter his teeth than milk.
}
{
	\aB{
		Nothing obvious has fallen out here. The hard word is חכלילי (\emph{ḥakhlili}), which is rare
		and obscure. This reconstruction connects it with an Akkadian word for darkness and takes the
		phrase as a construct/comparative: ``darkness of eyes more than wine''---in normal English,
		``his eyes are darker than wine.'' LXX has χαροποὶ οἱ ὀφθαλμοὶ αὐτοῦ ἀπὸ οἴνου (\emph{charopoi
		hoi ophthalmoi autou apo oinou}), an interpretive Greek rendering about his eyes and wine rather
		than a clean alternate Hebrew spelling. The second line is much easier: לבן שנים מחלב
		(\emph{lavan shinnayim meḥalav}) uses מן/מ־ in the normal comparative sense, so ``teeth whiter
		than milk,'' not ``teeth whitened by milk.'' The file again removes a line-opening \emph{waw}
		and writes the words in deliberately archaic spelling.
	}
}

\Verse{13}
{
	זבלן\textsuperscript{46}

	לחף ימים [ו]ישב

	על מפרצן\textsuperscript{47} ישכן

	\textsuperscript{48}\_\_\_\_\_\_

	\textsuperscript{49}\_\_\_\_\_\_
}
{
	Zebulun:

	At the shore of the seas he settles,

	At its inlets, he encamps.

	\ldots\ldots\ldots

	\ldots\ldots\ldots
}
{
	\aB{
		MT's Zebulun verse is readable but geographically and poetically suspicious: he lives by the
		seashore, is by a harbor of ships, and his ``flank'' reaches toward Sidon---even though
		historical Zebulun does not reach Sidon. The reconstruction borrows a much cleaner old parallel
		from Judg. 5:17, where the sea-shore/landing-place language survives, and rearranges it into
		``at the shore of the seas he settles / at its inlets he encamps.'' The endnote then tries to
		rebuild another pair of lines about sailing ships and vessels of Sidon from Phoenician and other
		Semitic parallels, but the file leaves those lines blank because the repair is too speculative.
		LXX just gives a smooth version of the MT idea---παράλιος κατοικήσει ... παρ᾽ ὅρμον πλοίων ...
		ἕως Σιδῶνος (\emph{paralios katoikēsei ... par' hormon ploiōn ... heōs Sidōnos})---so Greek does
		not rescue the lost wording here.
	}
}

\Verse{14}
{
	יששכר חמר \_\_\_\textsuperscript{50}

	רבץ בין [־]\textsuperscript{51} משפתם\textsuperscript{52}
}
{
	Issachar is an ass \ldots

	He couches among the rubbish piles
}
{
	\aB{
		Two old words are doing all the damage. MT calls Issachar חמר גרם (\emph{ḥamor garem}).
		``Donkey'' is clear; גרם is not. ``Bony/strong donkey'' is traditional, somebody proposed
		``castrated donkey,'' and the endnote offers a tiny emendation to גדים/גידים (\emph{gidim},
		``sinewy''). The file refuses to pick and leaves the adjective blank. Then comes משפתים
		(\emph{mishpetayim}), another rare old poetic word. It is often translated ``sheepfolds'' or
		``saddlebags,'' but this reconstruction connects it with אשפת (\emph{ashpot}), ``rubbish heap,''
		partly because related forms occur in other very old poems. LXX is radically different---Ἰσσαχαρ
		τὸ καλὸν ἐπεθύμησεν (\emph{Issachar to kalon epethymēsen}), ``Issachar desired the good''---so
		the Greek translator either had a different text or simply gave up on the donkey phrase too.
		Either way it cannot tell us the exact lost Hebrew.
	}
}

\Verse{15}
{
	ירא\textsuperscript{53} מנחה\textsuperscript{54} כי טוב

	יבט\textsuperscript{55} ארצו\textsuperscript{56} כי נעם

	יט\textsuperscript{57} שכמו לסבל

	ויה למס עבד\textsuperscript{58}
}
{
	He sees for himself a resting place
	which is good

	He beholds for himself a land
	which is pleasant

	He bends his shoulder to bear burdens

	And he has become a corvée slave.
}
{
	\aB{
		This verse is much less broken than the previous one. The reconstruction reads MT מנחה as מנחו
		(\emph{menuḥo}), ``his resting place,'' and הארץ as ארצו (\emph{artso}), ``his land.'' In really
		early spelling those final \emph{-o} vowels could be unmarked, and the Peshitta helps support
		the possessive readings. It also inserts יבט (\emph{yabbit}), ``he looks/beholds,'' before ``his
		land,'' because ``he sees / he looks'' is a standard poetic pair and the letters could have
		dropped beside a very similar sequence. Articles, the object marker, and a few opening
		\emph{waw}s are removed as later smoothing or metrical clutter. None of that changes the basic
		scene much: Issachar sees good land, puts his shoulder under the load, and ends up doing מס עובד
		(\emph{mas ʿoved}), compulsory state labor/corvée. That last phrase is not a textual problem; it
		is just an ugly social condition.
	}
}

\Verse{16}
{
	דן ידן <דין>\textsuperscript{59} עם

	כאחד שבט ישראל
}
{
	Dan pleads the cause of his people,

	The tribes of Israel together.
}
{
	\aB{
		MT already makes decent sense: דן ידין עמו (\emph{Dan yadin ʿammo}), ``Dan will judge his
		people.'' The reconstruction adds one repeated-looking word, דין (\emph{din},
		``case/judgment''), after ידין: ``Dan will judge the case of his people'' or ``plead his
		people's cause.'' The reason is meter plus an old Semitic idiom in which the verb ``judge''
		takes its matching noun; similar phrases occur in Jeremiah and Ugaritic. It would also be very
		easy for דין דין (\emph{din din}) to lose one copy through eye-skip. But no ancient manuscript
		actually preserves the extra noun here. This is a plausible little metrical restoration, not a
		necessary correction to an unintelligible MT.
	}
}

\Verse{17}
{
	[ ]\textsuperscript{60} דן נחש על דרך

	שפפן על ארח

	הנשך עקבי\textsuperscript{61} סס

	ויפל\textsuperscript{62} רכב\textsuperscript{63} אחר\textsuperscript{64} [ ]\textsuperscript{65}
}
{
	Dan is a serpent on the road,

	A venomous snake on the path

	Who snaps at the heels of a horse,

	So that its rider falls backward.
}
{
	\aB{
		MT begins יהי דן נחש (\emph{yehi Dan naḥash}), ``Let Dan be a snake.'' The file drops יהי as
		extra and treats this as a fresh little Dan saying. LXX actually keeps the verb---καὶ γενηθήτω
		Δαν ὄφις (\emph{kai genēthētō Dan ophis}), ``and let Dan become a snake''---so the deletion is
		meter-based, not manuscript-backed. Farther down, LXX has the simple singular picture δάκνων
		πτέρναν ἵππου ... πεσεῖται ὁ ἱππεὺς (\emph{daknōn pternan hippou ... peseitai ho hippeus}),
		``biting a horse's heel ... the rider falls,'' which helps explain the file's simplified old
		spellings, and the Peshitta supports the preferred reading of ויפל. Also: yes, MT verse 18 is
		missing from this file on purpose. MT has לישועתך קויתי יהוה (\emph{lishuʿatekha qiwwiti YHWH}),
		``For your salvation I wait, YHWH.'' The endnote treats that as a later worship line inserted
		around the midpoint of the tribal list. All the main ancient versions know the line, so deleting
		v. 18 is a literary-history judgment, not recovery of a manuscript that lacks it.
	}
}

\Verse{19}
{
	גד גדד\textsuperscript{66} יגדנו\textsuperscript{67}

	והא יגד עקב דם\textsuperscript{68}
}
{
	Gad is a band that raids

	And he raids from the rear.
}
{
	\aB{
		This is mostly one giant Gad pun: גד / גדוד / יגוד (\emph{Gad / gedud / yagud}), ``Gad /
		raiding-band / raid.'' MT normally gets something like ``a raiding band raids him, but he raids
		at their heels.'' The reconstruction reads יגדנו not as ``raids \emph{him}'' but as an old
		energic verb form, \emph{yegudun(na)}, with no object suffix: ``Gad is a band that raids.'' Then
		it takes the מ (\emph{m}) that MT attaches to the beginning of v. 20---מאשר (\emph{me-Asher}),
		``from Asher''---and moves it back to the end of this line with עקב (\emph{ʿaqev}), making an
		adverbial ``from the rear.'' That boundary move has real versional support: LXX, Vulgate, and
		Peshitta all begin the next saying simply with ``Asher,'' not ``from Asher.'' This is a nice
		example where one tiny letter probably got attached to the wrong verse rather than a whole
		sentence disappearing.
	}
}

\Verse{20}
{
	[ו־]\textsuperscript{69} אשר שמן לחמו\textsuperscript{70}

	שמנה ארצו

	[ ]\textsuperscript{71} יתן\textsuperscript{72} מעדני מלך
}
{
	As for Asher,
	his food is rich
	his land is fertile

	He produces royal delicacies.
}
{
	\aB{
		Once that stray מ is moved back into v. 19, MT still has a grammar problem: שמנה לחמו
		(\emph{shemena laḥmo}) puts a feminine ``fat/rich'' word next to masculine ``his bread/food.''
		The witnesses look like two old versions got mashed together. The Samaritan Pentateuch, LXX, and
		Vulgate support שמן לחמו (\emph{shamen laḥmo}), ``his food is rich''; LXX has Ἀσήρ, πίων αὐτοῦ ὁ
		ἄρτος (\emph{Asēr, piōn autou ho artos}), ``Asher, his bread is rich.'' The Peshitta and Targum
		seem to reflect another Hebrew line, שמנה ארצו (\emph{shemena artso}), ``his land is
		rich/fertile.'' The reconstruction's neat trick is to print \emph{both} lines and treat MT as a
		conflation of them. It then drops והוא (\emph{wehu}, ``and he'') as duplicated metrical padding. יתן (\emph{yitten}) could also be vocalized passively, ``he receives,'' but the file takes the
		active sense: he supplies royal delicacies.
	}
}

\Verse{21}
{
	\_\_\_\_\_\_

	\_\_\_\_\_\_\textsuperscript{73}
}
{
	\ldots\ldots\ldots

	\ldots\ldots\ldots
}
{
	\aB{
		MT gives the famous but weird נפתלי אילה שלחה הנתן אמרי שפר (\emph{Naftali ayyala sheluḥa,
		hannotein imre shefer}), usually ``Naphtali is a doe let loose, who gives beautiful
		words''---or, with other guesses, beautiful fawns. The LXX is startlingly different: Νεφθαλι
		στέλεχος ἀνειμένον, ἐπιδιδοὺς ἐν τῷ γενήματι κάλλος (\emph{Nephthali stelechos aneimenon,
		epididous en tō genēmati kallos}), roughly ``Naphtali is a spreading shoot/stem, producing
		beauty in its growth.'' That strongly suggests that very early readers had either a different
		Hebrew line or a radically different understanding of the same damaged consonants. The endnote
		tries to reconstruct אלה (\emph{ela}, ``terebinth'') and אמרי (\emph{amire}, ``foliage/tops''),
		making a tree that puts out beautiful leaves. But it openly admits the Greek words do not map
		neatly onto that Hebrew. That is why the file leaves the verse blank: LXX tells us
		\emph{something} is wrong, but not enough to recover exactly what was there.
	}
}

\Verse{22}
{
	\_\_\_\_\_\_

	\_\_\_\_\_\_\textsuperscript{74}

}
{
	\ldots\ldots\ldots

	\ldots\ldots\ldots
}
{
	\aB{
		This is where the Joseph poem really falls apart. MT has בן פרת יוסף בן פרת עלי עין בנות צעדה עלי שור (\emph{ben porat Yosef, ben porat ʿale ʿayin, banot tsaʿada ʿale shur}), traditionally
		turned into ``Joseph is a fruitful bough by a spring; branches climb over a wall.'' But פרת
		(\emph{porat}) is obscure, ``daughters marched'' is grammatically odd, and the Samaritan text
		has a different phrase around בני (\emph{beni}, ``my son''). The endnote experiments with
		reading the line as military language---something like a son marching against a spring and wall,
		natural targets in a siege---but does not claim victory. LXX is wildly different again: υἱὸς
		ηὐξημένος Ἰωσήφ ... υἱός μου νεώτατος· πρός με ἀνάστρεψον (\emph{huios ēuxēmenos Iōsēph ...
		huios mou neōtatos, pros me anastrepson}), roughly ``Joseph is a grown/flourishing son ... my
		youngest son, turn back to me.'' We have no Qumran fragment for this stretch. With witnesses
		this far apart, the file's blanks are the honest answer: the original wording cannot be
		reconstructed confidently.
	}
}

\Verse{23}
{
	\_\_\_\_\_\_

	\_\_\_\_\_\_\textsuperscript{75}
}
{
	\ldots\ldots\ldots

	\ldots\ldots\ldots
}
{
	\aB{
		Ironically, this verse by itself is not nearly as hopeless as the blank makes it look. MT
		describes hostile archers attacking Joseph, and בעלי חצים (\emph{baʿale ḥitstsim}), literally
		``masters of arrows,'' is odd but perfectly understandable as ``archers.'' The main preferred
		correction comes from the Samaritan Pentateuch: יריבהו (\emph{yerivuhu}), ``they strove against
		him,'' instead of MT's difficult ורבו (\emph{warobbu}). That improves both sense and meter. LXX
		keeps the same general attack scene: εἰς ὃν διαβουλευόμενοι ἐλοιδόρουν ... κύριοι τοξευμάτων
		(\emph{eis hon diabouleuomenoi eloidoroun ... kyrioi toxeurmatōn}), people plotting/taunting him
		and ``lords of arrows.'' The file blanks the verse because vv. 22--24 form one damaged unit and
		nobody can confidently reconnect all the pieces, not because every word in v. 23 is unreadable.
	}
}

\Verse{24a}
{
	\_\_\_\_\_\_

	\_\_\_\_\_\_\textsuperscript{76}
}
{
	\ldots\ldots\ldots

	\ldots\ldots\ldots
}
{
	\aB{
		MT can be translated into something smooth---``his bow remained firm, and the arms of his hands
		were agile/strong''---but the Hebrew underneath is nasty enough that the endnote gives up. The
		LXX does not stabilize it; it gives almost the opposite picture: καὶ συνετρίβη μετὰ κράτους τὰ
		τόξα αὐτῶν, καὶ ἐξελύθη τὰ νεῦρα βραχιόνων χειρῶν αὐτῶν (\emph{kai synetribē meta kratous ta
		toxa autōn ...}), ``their bows were broken with force, and the sinews of their arms were
		loosened.'' Notice even the possessors flip from Joseph's bow/arms in MT to \emph{their}
		bows/arms in Greek. That could reflect a different Hebrew text, ancient interpretation, or
		corruption on either side. There is no clean way back, so the file leaves both lines blank.
	}
}

\Verse{24b}
{
	\hProto{
		מיד אביר יעקב

		משמר\textsuperscript{77} ר־בני ישראל
		מרע
	}
}
{
	\eProto{
		From the hands of the Champion of Jacob

		From the Keeper

		From the Shepherd of the sons of Israel
	}
}
{
	\aB{
		MT's ending is מידי אביר יעקב משם רעה אבן ישראל (\emph{mide avir Yaʿaqov, misham roʿe even
		Yisraʾel}), literally something like ``from the hands of the Mighty One of Jacob; from there,
		shepherd, stone of Israel.'' ``Stone of Israel'' can be theologized into a divine title, but the
		line is long and the syntax is rough. This reconstruction thinks משם רעה is actually two old
		readings accidentally fused together: משמר (\emph{mishshomer}), ``from the Keeper,'' and מרעה
		(\emph{meroʿe}), ``from the Shepherd.'' It then changes אבן (\emph{even}, ``stone'') to בני
		(\emph{bene}, ``sons of''), perhaps with the stray aleph attracted by nearby אביר (\emph{avir}).
		That yields a tidy set of titles: ``from the Champion of Jacob / from the Keeper / from the
		Shepherd of the sons of Israel.'' Ps. 121:4 calls YHWH the keeper of Israel and Ps. 80:2 the
		shepherd of Israel, so the language is real even if this exact reconstruction is conjectural.
		LXX has yet another line about the mighty one of Jacob and ``the one who strengthened Israel,''
		so Greek is not the direct source of the file's repair.
	}
}

\Verse{25}
{
	\hProto{
		מאל אבך\textsuperscript{78} ויעזרך

		ואל־שדי\textsuperscript{79} ויברכך

		ברכת שמם מעל\textsuperscript{79a}

		ברכת תהם\textsuperscript{80} [ר] מתחת

		ברכת שדם ורחם
	}
}
{
	\eProto{
		From the God of thy father,
		who supports thee,

		And El-Shaddai, who blesses thee:

		Blessings of heaven above,

		Blessings of the Deep, beneath,

		Blessings of breasts and womb,
	}
}
{
	\aB{
		The first repair has actual manuscript support. MT has ואת שדי (\emph{we-et Shaddai}) in the difficult sequence around Shaddai, while three
		Hebrew manuscripts, the Samaritan Pentateuch, and Peshitta
		support ואל שדי (\emph{we-El Shaddai}), ``and El-Shaddai,'' which the file adopts. LXX instead
		has ὁ θεὸς ὁ ἐμός (\emph{ho theos ho emos}), ``my God,'' so there was already variation in this
		little divine-title line. The next problem is the blessing of the Deep. MT says ברכת תהום רבצת תחת (\emph{birkhot tehom rovetset taḥat}), ``blessings of the Deep crouching beneath.'' The
		reconstruction drops רבצת (\emph{rovetset}, ``crouching'') because the line is too long and
		looks harmonized toward the parallel Joseph blessing in Deut. 33. LXX goes farther and has καὶ
		εὐλογίαν γῆς ἐχούσης πάντα (\emph{kai eulogian gēs echousēs panta}), ``and a blessing of earth
		containing everything,'' close in idea to Deut. 33:16. The important beginner point is that the
		Joseph blessing seems to have circulated in more than one form; Genesis 49 and Deuteronomy 33
		can sometimes correct each other, but neither is simply a photocopy of one lost master text.
	}
}

\Verse{26}
{
	\hProto{
		ברכת אב <ואם>\textsuperscript{81}
		גבר ועל

		ברכת הרי\textsuperscript{82} עד

		ברכת\textsuperscript{83} גבעת עלם

		תהיין\textsuperscript{84} לראש יסף

		ולקדקד נזר\textsuperscript{85} אחיו\textsuperscript{86}
	}
}
{
	\eProto{
		Blessings of
		father and mother
		man and child

		Blessings of the mountains of old,

		Blessings of the eternal hills;

		Let them be on the head of Joseph

		On the brow of the leader
		of his brethren.
	}
}
{
	\aB{
		MT's opening is badly tangled: ברכת אביך גברו על ... (\emph{birkkhot avikha gaveru ʿal ...}),
		then a phrase about ancient heights/hills that is hard to parse. The Samaritan Pentateuch and
		LXX preserve a much clearer pair, ``blessings of father and mother''; Greek explicitly has
		εὐλογίας πατρός σου καὶ μητρός σου (\emph{eulogias patros sou kai mētros sou}). The endnote
		thinks MT has conflated that with another old pair about ``man and child,'' so the file prints
		both ideas instead of forcing the mashed line into one sentence. Next, MT הורי (\emph{horai}) is
		read as הרי (\emph{harare/hare}), ``mountains,'' helped by the same archaic phrase in Hab. 3:6
		and by Deut. 33. Then MT תאוה (\emph{taʾawat}, ``desire'') is changed to ברכות (\emph{birkhot},
		``blessings''); here LXX gives direct support with εὐλογίαις (\emph{eulogiais}), ``blessings,''
		and the word is already repeating all through the poem. The closing ``on Joseph's head / on the
		brow of the leader among his brothers'' is much safer. So the middle of v. 26 is a classic
		conflated, cross-contaminated old poem; the end is basically intact.
	}
}

\Verse{27}
{
	\hProto{
		בנימן זאב יטרף\textsuperscript{87}

		בבקר יאכל עד

		לערב\textsuperscript{88} יחלק שלל\textsuperscript{89}
	}
}
{
	\eProto{
		Benjamin is a wolf who preys,

		In the morning he devours spoil,

		At evening he divides booty.
	}
}
{
	\aB{
		This is another clean ending. MT's בנימין זאב יטרף (\emph{Binyamin zeʾev yitraf}) literally has
		``Benjamin, a wolf preys'' without an אשר (\emph{asher}, ``who/that''); that missing relative
		word is normal in very old Hebrew poetry, not damage. The file therefore reads naturally,
		``Benjamin is a wolf who preys.'' It drops the \emph{waw} before the evening line because
		line-opening ``and'' is often absent in early verse. LXX softens the last image into εἰς τὸ
		ἑσπέρας διαδώσει τροφήν (\emph{eis to hesperas diadōsei trophēn}), ``at evening he distributes
		food,'' while MT has יחלק שלל (\emph{yeḥalleq shalal}), ``he divides spoil.'' That looks more
		like translation/interpretation than evidence for a missing Hebrew line. Basically, after all
		the wreckage in Joseph, Benjamin lands us back on solid textual ground.
	}
}

\section*{Notes to the Blessing of Jacob}

\begin{enumerate}

\item
Following Prof. Albright's suggestion, we reconstruct vss. 1b--2 as a
tricolon, 3:3:3. The words in vs. 1b beginning with \emph{'ēt 'ăšer} and
continuing to the end of the verses seem to be a prosaic addition. In addition,
\emph{wəšim'û} after \emph{hiqqābĕṣû} is to be omitted, apparently having been
inserted under the influence of \emph{wəšim'û} in the second colon.

\item
The pronoun is to be connected with the following phrase, in agreement with
the LXX and Vulgate. On the meter, see note 4.

\item
Omit the conjunction, following several LXX minuscules. According to the
canons of early Canaanite and Hebrew poetry, the conjunction is used rarely at
the beginning of cola.

\item
The difficult phrase \emph{yeter śə'ēt} appears to be a corruption of
\emph{rē'šît}, repeated accidentally from the previous colon. The original
reading of the final colon of the verse may have been either \emph{rē'šît
'uzzî} (on \emph{'uzzî}, see note 5), or \emph{yeter 'uzzî}. Since the present
text may be more easily explained on the basis of an original \emph{yeter},
and since that reading is supported by the versions, it is to be preferred.

Analysis indicates that the metrical structure of vs. 3 is 2:2:2:2. Note the
rhyme scheme, typical of bedouin poetry, in which each colon ends with the
first person pronominal suffix, \emph{-î}.

\item
Read \emph{'uzzî}, ``my strength.'' The suffix is to be read, as indicated by
the parallel expressions: \emph{bəkōrî, kōḥî}, and \emph{'ōnî}. In עז, we
have an example of early Hebrew orthography, final vowels not being indicated
in the spelling. Ball, \emph{The Book of Genesis (Sacred Books of the Old
Testament)}, Leipzig, 1896, p. 119, suggests the reading \emph{'uzzî}. In the
parallel phrase, he reads \emph{yeter śə'ētî}.

\item
Compare with this couplet the Song of Lamech, Gen. 4:23--24, in which we find
a similar rhyme and a similar metrical pattern.

\item
Vs. 4a is unintelligible as it stands. The versions are based on an already
corrupt text and offer no help in interpreting the passage. The line apparently
is defective, a bicolon 3:3 being expected here (since this material goes with
what follows). Possibly the first two words are to be read, ``thou art
uncontrollable, like the (waters of the) sea.'' The next phrase is undoubtedly
corrupt; cp. the similar constructions in vs. 6 with \emph{'al}.

\item
The verb \emph{'ālā} is used customarily with the preposition \emph{'al} in a
construction of this type; cf. Ps. 132:3. It has been lost, apparently as the
result of homoiarkton, and is to be supplied after the verb (the long form
\emph{'ālē} suits the metrical requirements). The anomalous \emph{'ālā} at the
end of the verse may reflect a later restoration of the missing preposition.

\item
Read \emph{miškābî}, ``my bed.'' The word is probably singular, since the
plural form, with this meaning, is always \emph{miškābôt}. The apparent
exception, \emph{miškəbê} (Lev. 18:22; 20:13), is in reality a biform with a
different meaning. On the displacement of \emph{'ābîkā}, see note 11.

\item
Read either the construct plural or the construct singular with the genitive
case ending, following the Massoretic text. On the use of the \emph{litterae
compaginis}, see GK §90l, and Albright, \emph{JBL} LXI (1942), p. 117; and
``The Oracles of Balaam,'' \emph{JBL} LXIII (1944), p. 215, note 48; p. 216,
note 54; etc.

\item
On the shift of \emph{'ābîkā} from the end of the first colon to the end of
the second colon, see the parallel text in I Chron. 5:1, which is dependent
upon this verse. Note that this rearrangement of the text clears up the
obscurities in this part of the verse, and that the resulting meter is a quite
regular 3:3.

\item
There has been considerable discussion as to the precise significance of
\emph{'aḥîm} in this context. The usual interpretations emphasize either the
genetic relationship (as uterine brothers) or the spiritual kinship (as
partners in crime). The underlying thought probably is historical: there was a
close association between the tribes historically, and in the blessing they
are regarded as having a common origin and a common fate (compare the account
of their joint exploits against the Shechemites in Gen. 34).

\item
This word is a \emph{hapax legomenon}. A number of different etymologies have
been suggested, but as yet no satisfactory solution has been achieved. The
word seems to fall into one of two possible areas of meaning: (1) it refers to
some kind of weapon, possibly battle-axe, and is to be derived from the root
\emph{kwr, krr}, or \emph{krt}; (2) it is connected with the root \emph{mkr},
and has the meaning ``merchandise,'' and in this context, ``stock in trade,''
i.e., ``implements of violence are their stock in trade.'' In this connection,
note Hebrew \emph{makkār} (II Kings 12:6, 8), ``broker'' (so Albright), and
Ugaritic \emph{mkr} with the same meaning.

The final suffix probably is to be read \emph{-ām}, with Sievers and Ball.

\item
It is preferable to translate this colon as direct address: ``Into their
council enter not, O my soul! In their assembly, do not join, O my heart
(lit. liver, seat of the emotions)!'' Read, therefore, \emph{tābō'î}, second
feminine singular; in early orthography the final yodh would not be written.
This reading is indicated by \emph{tēḥad} in vs. 6a, which must be construed
as a second masculine singular. For other examples of \emph{napšî} used in
direct address, see Pss. 103:1 and 104:1. Cf. Peters, ``Jacob's Blessing,''
\emph{JBL} VI (1886), p. 101.

\item
The expression here is elliptical; \emph{'ittām}, ``with them,'' is to be
understood after \emph{tēḥad}. The complete construction is to be found in
Isa. 14:21. The Samaritan text, יחד, represents a revision of the Hebrew based
on a misunderstanding of the syntax of the sentence; since the subject was
masculine, and the sentence was read in the third person, the verb was revised
to agree with the noun. Subsequently, \emph{resh} was substituted for
\emph{daleth}, a common error at all periods of scribal transmission. The LXX
reads μὴ ἐρίσαιτο τὰ ἥπατά μου, ``let not my emotions be aroused.'' This
interpretation is based upon the same text as the Samaritan, \emph{yiḥar},
with the additional change to \emph{kəbēdî}, ``my liver''; cf. note 16.

\item
Read probably \emph{kəbēdî}, following the LXX. MT may be correct; cf. König,
\emph{Die Genesis}, Gütersloh, 1919, \emph{ad loc.}; but the Ugaritic usage
points to the other reading; \emph{kbd}, ``liver,'' occurs in parallel with
\emph{lb}, ``heart,'' and \emph{'irt}, ``chest, lungs.'' This suggestion comes
from Prof. Albright.

\item
Omit \emph{kî}, which, as it stands, impairs the symmetry of the two cola.

\item
Omit the conjunction; cf. note 3.

\item
There is no reason to look for either historical or mythological reflections
in this couplet. The device used here is common to early Canaanite and Hebrew
poetry. By the employment of non-exact, highly figurative language, a striking
and dramatic picture is created. This is impressionistic poetry at its best.
What is involved here is not a reference to the particular slaying of men or
the maiming of beasts, but a general though sharp impression of the lawlessness
and violence of the two tribes. A similar idea is conveyed by the description
in the Song of Lamech, Gen. 4:23--24. Other striking examples occur in
Deut. 33:18 and Jud. 5:25--26. Examples from the Ugaritic literature are to be
found in Gordon's \emph{Ugaritic Handbook}, I, pp. 102ff.

Zimmern, ``Der Jakobssegen und der Tierkreis,'' \emph{ZA} VII (1892),
pp. 161--172, finds in this verse a mythological borrowing from the Babylonian
Gilgamesh. Gilgamesh and Enkidu are recorded as slaying a giant named Huwawa =
Humbaba, the guardian of the cedar forests; and in another episode, as together
killing the Bull of heaven. Following this, Ishtar pronounces a curse on
Gilgamesh, and Enkidu is slain by the gods. Certainly no direct connection
between the Babylonian myth and the Israelite poem exists. If there is any
relationship at all, it is a remote one, and lies in the use of poetic language
in the blessing which may ultimately rest upon an allusion to the heroic
exploits of the Babylonian brothers-in-arms. To the composer of the blessing,
the expression was a striking description of a pair of violent men. That is all
it means in this context. It may possibly, however, have come down to him as a
brief summary of the activities of those typical men of violence, Gilgamesh
and Enkidu. See also E. Burrows, \emph{The Oracles of Jacob and Balaam},
London, 1939; his conclusions are utterly fantastic, however. Cf. note 26.

\item
The use of the first person in this blessing presents no real difficulty in
interpretation. The blessing is put in the mouth of the eponymous ancestor,
Jacob-Israel. On vs. 7b, see Gunkel, \emph{Genesis}, \emph{ad loc.},
``Jacob selber zerstreut sie, eben durch diesen Fluch.'' Also see König,
\emph{Genesis}, \emph{ad loc.}

\item
In this blessing the metrical balance between parallel cola is strictly
observed. It is more difficult to determine the accentual pattern. Apparently
the pattern is 3:3:3:3 in vss. 5--6a; for 6b--7, a 2:2:2:2:2:2 structure
emerges. The dividing line is uncertain, however, and some of these cola can
be construed as having three accents. This, however, does not affect the
balance between cola, where in almost every case we have an equal number of
syllables.

\item
Restore \emph{hôdî}, ``my majesty,'' or the like. A word has fallen out of the
first colon, and this is the simplest emendation of the text. Written in the
old orthography, \emph{yəhûdâ} would have appeared as יהד; \emph{hôdî} as הד.
The loss of the second word would result from haplography.

\item
Contrary to the commentaries, \emph{'attâ} can hardly be construed in the
accusative case, in apposition with the verbal suffix, \emph{-kā}. Examples in
GK §135e of this construction are unsatisfactory, being either the result of
textual corruption or not strictly applicable to the present case. According
to the present Massoretic text, we must read ``Judah (art) thou; thy brothers
praise thee!'' On the first part of the colon, see note 22. On the structure
of the blessing, see note 24.

\item
The clue to the structure of this couplet is to be found in the parallelism
between \emph{yādəkā} and \emph{yištaḥăwû}, and \emph{'aḥêkā} and
\emph{bənê 'ābîkā}. As the text now stands, the colon \emph{yādəkā bə'ōrep
'ōyəbêkā} is awkwardly placed between the two cola referring to Judah's
brothers; however, it fits well with the reconstructed first colon of the
blessing. On the expression \emph{bə'ōrep 'ōyəbêkā}, cf. Ps. 18:41 =
II Sam. 22:41, and Ex. 23:27.

\item
This colon has been reconstructed on the basis of the LXX, Vulgate, and
Peshitta, none of which read the independent pronoun, but only the pronominal
suffix; cf. note 22; on the shift of the colon from the beginning of the verse,
see note 24. Apparently very early in the textual transmission of this passage,
the first and third cola were confused, the present MT preserving parts of
both mixed together, while the text from which the versions were made
preserved only the third colon, shifted, however, to the beginning of the
blessing. It is to be noted that the name ``Judah'' is repeated a number of
times through the blessing, vss. 9 and 10, as well as in vs

% ---- The note above is 25. It is followed in the text by 25a.

% ---- The format abruptly changes here in the OCR output. Make it consistent.

25a. Cf. Isaac's blessing of Jacob in Gen. 27:29.

26. Zimmern's speculations about the different tribes and their zodiacal
associations are ingenious but entirely unacceptable. It is quite certain that
the zodiac did not appear in developed form until the Persian period. Even if
a few of the elements which were later incorporated into the zodiacal system
were known in earlier times, none of the parallels between the signs of the
zodiac and the designations of the tribes is at all convincing. Cf. Ball,
\emph{Genesis}, pp. 114ff. On the designation of Judah here, note that Dan is
described as a lion's whelp in Deut. 33:22.

27. The precise meaning of this colon is not certain. The usual explanation,
i.e. that the lion rises from his prey and ascends to his mountain lair, seems
the best so far suggested.

28. This old poetic phrase was also applied to the nation as a whole in the
\emph{Oracles of Balaam}, Num. 24:9 (cf. Albright, ``The Oracles of Balaam,''
\emph{ad loc.}).

29. Omit the conjunction at the beginning of the second colon, \emph{metri
causa}.

30. On this construction, see Chapter IV, note 35.

31. Read \emph{šōpēṭ}, ``judge, charismatic leader,'' or better perhaps,
following Albright, \emph{šēpēṭ}, Ugar. \emph{ṭpṭ}, with the same meaning; cp.
the LXX, ἄρχων. In most instances, \emph{məḥōqēq} means ``commander''; cf.
note 32. In Jud. 5:14, where it clearly has this meaning, it stands in parallel,
not with \emph{šēbēṭ}, but with the ``wielder of the \emph{šēbēṭ}.'' In
Isa. 32:22, \emph{šōpēṭ} and \emph{məḥōqēq} are used synonymously, a fact which
strongly supports the present emendation. In the early script, \emph{pe} and
\emph{beth} were very much alike in form, and thus easily confused; by the same
token they sound very much alike and the confusion may arise in oral
transmission. Note that in Ps. 68:5, we have \emph{ʿarābôt}, where the
\emph{beth} is an error, the correct original being \emph{pe}. A striking
instance of the confusion between \emph{šēbēṭ} and \emph{šōpēṭ} (or
*\emph{šēpēṭ}) is to be found in the parallel texts, II Sam. 7:7 =
I Chron. 17:6, where II Samuel reads, incorrectly, \emph{šibṭê}, while
I Chronicles preserves the correct original, \emph{šōpəṭê}.

32. The meaning ``commander'' is certain for \emph{məḥōqēq} in the following
passages, Deut. 33:21; Jud. 5:14 (cp. Jud. 5:9, \emph{ḥōqəqê}, which is an
abbreviation for \emph{məḥōqəqê}, either deliberate or accidental); and
Isa. 33:22. In two other instances, the meaning is not so clear, Num. 21:18 and
Ps. 60:9 = Ps. 108:9. Morphologically, the word should mean, ``one who issues a
\emph{ḥōq},'' i.e., a commander.

33. Read with Samaritan, \emph{dəgālâw}, ``his standards,'' ``his battle
flags.'' The whole phrase would read, ``from the midst of his banners,'' an
excellent parallel to \emph{mîhûdâ} in the first colon. The expression
\emph{mibbên raglayim} is a euphemism referring to the female genitals; cf.
Deut. 28:57 (for other euphemistic usages, see lexicons). The phrase in MT
could only refer to the male counterparts, and indeed is so taken by the LXX
and Vulgate.

34. No satisfactory solution of the problems in this bicolon has ever been
presented, and there is no present prospect of a definitive solution. For a
brief summary of the situation, see Skinner, \emph{A Critical and Exegetical
Commentary on Genesis (International Critical Commentary)}, New York, 1910,
pp. 519--524. G. R. Driver, ``Some Hebrew Roots and Their Meanings,'' \emph{JTS}
23 (1922), pp. 69--73, suggests for the critical word the vocalization
שִׁילֹה, ``prince'' = Acc. \emph{šēlu, šīlu}. The following reconstruction
of the text may suggest the lines along which the problem should be attacked:

\begin{quote}
עדיו (עדך) יבא שין[ ] ולו יקהת עמים

\emph{ʿadêw} (var. \emph{ʿadêkā}) \emph{yûbaʾ šay wəlô yiqqəhat ʿammîm},

``Unto him (unto thee) shall tribute be brought, and to him shall be the
subservience of peoples.''
\end{quote}

35. Read \emph{ʾašûrî}, Qal passive participle, construct, with the
\emph{hireq compaginis}; cp. Samaritan אסורי.

36. Read probably \emph{ləgapnô}, ``to its harness.'' The suffix of the third
person masculine singular would not be indicated in the early orthography.
The meaning ``harness, trappings,'' for *\emph{gapn(u)} is established from
the Ugaritic texts; cf. 51:IV:9--12 (= 51:IV:4--7),
\emph{mdl. ʿr. ṣmd. pḥl. št. gpnm. dt. ksp. dt. yrq. nqbm. ʿdb. gpn.
ʾatnth}. The use of the preposition ל in the phrase \emph{ʾašûrî ləgapnô} is
not clear; if the interpretation of the passage suggested above is correct,
then it must be construed as a ל of reference or agency. Cf. GK §119u, §121f;
BDB, pp. 513--514.

For another interpretation of the passage which has received too little
attention from scholars, see Jastrow, ``Light Thrown on Some Biblical Passages
by Talmudic Usage,'' \emph{JBL} XI (1892), pp. 126ff. Also see A. Jirku,
``Der Juda-Spruch, Genesis 49:18ff, und die Texte von Ras Samra,'' \emph{JPOS}
XV (1935), pp. 12--13.

37. Read \emph{ʿērô} < \emph{ʿayrô}, Ugar. \emph{ʿr}. The preservation of
the \emph{he} to indicate the third masculine singular suffix reflects
9th--6th century B.C. spelling practice.

38. Omit the conjunction at the beginning of the second colon.

39. There is no evidence for a feminine form of this word, the other
occurrences of this word in the Bible being masculine. The final \emph{he} is
to be regarded therefore as the third masculine singular suffix (agreeing with
the masculine \emph{bənî ʾatōnô}); on the orthography, see note 37.

40. Note the preservation of the archaic case ending in the \emph{nomen
regens} of a construct chain, a not uncommon feature of early Hebrew poetry.
On the bicolon as a whole, compare Zech. 9:9.

41. Omit the \emph{waw}, possibly a case of dittography.

42. The word means clearly ``mantle, garment,'' and is found in the Kilamuwa
Inscription from the late ninth century, line 8, סות (the \emph{waw} is
consonantal), and in the Batnoʿam Inscription. Cf. P. Joüon, ``Gen. 49:11
סות,'' \emph{Biblica} 21 (1940), p. 58, among others on סות and its
vocalization. Samaritan כסותה seems to be an ancient emendation. On the final
\emph{he}, see note 37.

43. This is apparently a construct chain (cf. note 40), to be rendered
literally, ``darkness of eyes more than wine''; cp. the LXX paraphrase, ``his
eyes are darker than wine.'' For the meaning of \emph{ḥaklîlî}, note Acc.
\emph{eklu, eklītu}, ``dark, darkness.'' Cp. Prov. 23:29, where חכלילות means
``darkness,'' but with the specific sense of ``bloodshot.'' The Samaritan
וחכלילו is either a misspelling (\emph{waw} and \emph{yodh} are almost
identical in the Herodian script), or it may preserve the original and correct
nominative case ending, \emph{-u}.

44. Omit the \emph{waw} at the beginning of the second colon.

45. This expression shows that the preposition is used with the adjective in a
comparative sense, i.e., ``teeth whiter than milk,'' rather than ``teeth made
white by milk.'' So Jastrow and others.

The description in vs. 11 may well have a Canaanite mythical background; cf.
Sellin, ``Zu dem Judasspruch im Jakobssegen Gen. 49:8--12 und im Mosesegen
Deut. 33:7,'' \emph{ZAW} 60 (1944), pp. 57--67, esp. p. 63. Sellin's attempts
to date the Judah blessings in Gen. 49 and Deut. 33 are, however, most
unfortunate.

46. Both in MT and in the reconstructed text presented here, the word
\emph{zəbûlûn} stands as extra-metrical. This is not to be regarded as a
hypothetical occurrence in Hebrew poetry, since such extra-metrical feet are a
common device in the Ugaritic epics; cf. Ginsberg, \emph{Orientalia} 5 (1936),
pp. 171ff. In Deut. 33, in two instances, the name of the tribe is not included
in the blessing itself (Benjamin, vs. 12, and Joseph, vs. 13). Presumably the
name stood at the beginning of the blessing (disregarding the titles which now
stand in MT, a later development in the transmission of the blessings), as an
extra-metrical heading. In time this practice was generalized, and headings
were attached to the other blessings, whether the tribe was named in the body
of the blessing or not.

47. This bicolon is reconstructed on the basis of Jud. 5:17, where the complete
text is preserved. By a slight rearrangement of the words (i.e., \emph{yāšab}
after \emph{yammîm}) the omission in Gen. 49 becomes explicable as a case of
homoiarkton. The description was applicable to the whole area in which the
tribes of Asher and Zebulun were located, and could therefore be used with
equal facility for either.

48. The word \emph{ḥōp} does not belong in this context. In every other
instance in the Bible, it is coupled with the word for ``sea'' (as in vs.
13a); its insertion here is the result of dittography. It has been suggested
by a number of scholars that לחף is a corruption of חבל, and that is quite
possible. Equally possible, and perhaps preferable, is the reading לפף,
suggested by Prof. Albright. The colon would then read ``And he indeed doth
fare on ships.'' Cf. לפף שבל, at the end of the Ahiram Inscription, which is
to be translated ``wayfarer''; Albright, ``The Phoenician Inscriptions of the
Tenth Century B.C. from Byblus,'' p. 155.

49. As it stands, this bicolon makes no satisfactory sense. There is no
parallelism between the cola, and the geographical reference is wrong on
topographical grounds (i.e., Zebulun did not extend to Sidon, so על or עד is
incorrect). If the first colon is restored in accordance with the suggestion
above (note 48): \emph{wəhû' lōtēn 'ŏniyyôt}, then some equivalent idea must
be present in the second colon. On the basis of Isa. 33:21 (cp. Num. 24:24),
צי is restored before ציד; its loss being due to simple haplography. In this
context, ירכב is impossible; the least drastic emendation would be ירכב. The
resulting text is: \emph{yirkab 'al ṣiyyê Ṣîdōn}. The use of \emph{rkb} for
travelling on ships is frequent in Accadian, and regular in Arabic. The full
bicolon would then be translated.

\begin{quote}
And he is a sailor of ships\\
He rides on the vessels of Sidon.
\end{quote}

50. This word remains obscure. The usual interpretation, ``bony,'' i.e.,
``strong,'' is unsatisfactory. The best recent treatment of the phrase
\emph{ḥămōr gārēm} is that of Feigin, ``Hamor Garim, `Castrated Ass,' ''
\emph{JNES} V (1946), pp. 230--233. It is possible, however, to emend the text
slightly, and read \emph{ḥămōr gidîm}, ``sinewy ass''; this makes better sense
in the context.

51. Omit the article, which is a later insertion. It is not found in early
Hebrew or Canaanite poetry.

52. This word occurs in three early passages of Hebrew poetry: Jud. 5:16,
Ps. 68:14 (without the \emph{mem} preformative), and here. It is presumably to
be associated with \emph{'ašpōt}, ``rubbish heap.'' On the etymology of the
term, cf. Albright, ``The Earliest Forms of Hebrew Verse,'' \emph{JPOS} II
(1922), p. 78, note 2.

53. Omit the conjunction at the beginning of the cola, as in the Vulgate.
Read \emph{yir'ê}, ``he saw.'' In the early orthography, the final \emph{he}
would not have been written.

54. Read \emph{mənûḥô}, ``his resting place,'' following a number of scholars.
The third masculine singular suffix was indicated by \emph{he} in pre-exilic
times; cf. note 37. This is implied by the masculine adjective \emph{ṭôb}. The
Peshitta also reads the suffix here; cf. note 56. The Samaritan offers a
variant text, מנוחה and טובה; the LXX and Vulgate have the same interpretation
of the text.

55. The \emph{nota accusativi} and article reflect a later revision of the
text; these may have been introduced to balance the meter in a defective line,
since something has apparently dropped out. It is suggested that
\emph{yabbîṭ}, ``he looked upon,'' has fallen out by haplography. Note that the
three preceding letters (in pre-exilic orthography) are \emph{yodh, ṭeth}, and
\emph{beth}. It is unnecessary to point out that \emph{rā'â} and \emph{hibbîṭ}
are used consistently in parallel construction throughout the Bible.

56. Read \emph{'arṣô}, ``his land,'' following the Syriac; cf. note 54. The
final vowel would not be indicated in tenth-century orthography.

57. Omit the conjunction (cf. LXX). Read \emph{yiṭṭê}; the \emph{he} would not
have appeared in early orthography.

58. On the expression \emph{mas 'ōbēd}, see the thorough treatment by
I. Mendelsohn, ``State Slavery in Ancient Palestine,'' \emph{BASOR} \#85
(1942), pp. 14--17, and \emph{Slavery in the Ancient Near East}, New York,
1949. It refers to a variety of state slavery in which the corvée workers were
reduced to the status of slaves. This was not practiced in Israel until the
time of David and Solomon, and then only upon subjugated non-Israelites. The
instance here must refer to the subjection of Issachar to the Canaanites at an
early period of the occupation, or perhaps to the Philistines toward the end of
the period of the Judges. Conceivably this may reflect the situation of certain
Hebrew groups in Palestine before the conquest.

59. This colon is considerably shorter than the following one; on the analogy
of vs. 19, and based upon a common construction with the verb \emph{dîn}, we
supply the noun \emph{dîn} after \emph{yādîn} (lost by haplography). The
passage then reads, ``Dan will plead the cause of his people.'' Compare the
following passages: Jer. 5:28; 22:16; 30:13; in Ugaritic, 127:33--34 =
127:45--47,

\begin{quote}
\emph{la-tadinu \qquad din(a) \qquad 'almanati}\\
\emph{la-taṭpuṭu \qquad ṭipṭa \qquad qaṣir(i) \qquad napši.}
\end{quote}

60. Vs. 17 constitutes a separate blessing, not directly connected with the
previous one. Omit יה, which is superfluous, disturbing the meter and the
meaning.

61. The LXX apparently read \emph{'āqēb} (singular) here; cp. also the
Peshitta, which paraphrases. This may reflect the older spelling practice, in
which the final vowel letters were not written; or it may be a correction on
the basis of a literal interpretation of the passage (i.e., the snake bites
only one heel), whereas the real meaning is less prosaic: ``Who snaps at the
heels of the horse. . . .'' The phrase \emph{'iqbê sûs} occurs also in the
Song of Deborah, Jud. 5:22.

62. Read \emph{way-yappēl}, following Ball (on the basis of the Peshitta).

63. The LXX apparently read \emph{rōkēb}; this may reflect the older
orthography, in which the final vowel would not be indicated. The clear
implication of the passage is that the horse, attacked by the snake, bolts and
throws its rider. This scene may reflect any period in the history of Palestine
after the Hyksos (18th century B.C.), who introduced horses and chariots in
considerable numbers. It is true that cavalry was not used in Palestine until
the ninth century, but there was a long period of informal or non-military
horseback riding before that. A good illustration of horseback riding in the
Amarna Age is to be found in Letter No. 245, in which Biridiya of Megiddo tells
how his mare was shot by an arrow and felled. Cf. Knudtzon, \emph{El-Amarna
Tafeln}, I, Leipzig, 1915, pp. 292--293, Biridiya of Megiddo to the King \#4
(Prof. Albright).

64. The metrical structure of the couplet (vs. 17) presents a problem. It
consists of a pair of bicola, which may be construed as either 2:2.2:2 or
3:3.3:3. In the former case, the first three cola present no particular
difficulty, but the last one clearly has three stresses; in the latter case,
the first, third, and fourth cola are satisfactory, but the second has only two
stresses. Nevertheless, the requirements of symmetry are amply met, and an
almost perfect balance obtains in the couplet as it stands. Note that there is
practically no variation in the number of syllables in each of the four cola.

65. This verse does not belong with the Blessing of Dan, nor is it an original
part of the blessings as a group. It seems to be a liturgical rubric inserted
approximately at the midpoint in the series of blessings, when the poem was
adapted for use in worship.

66. The word \emph{gədûd} is to be construed as a predicate nominative with
\emph{gād}. Contrary to the commentators, the play on words must turn on the
characterization of Gad, not his enemies. This is confirmed by the other
blessings in the poem, in which the tribes are characterized under different
figures: Judah, vs. 9, Issachar, vs. 14, Dan, vs. 17, Naphtali, vs. 21, and
Benjamin, vs. 27.

67. Read \emph{yəgûdun(na)}, the energic form without the suffix. This is a
common construction in Ugaritic, and occurs a number of times in early
Israelite poetry as well; cf. Chapter IV, note 25 for examples from the
Blessing of Moses and references to the Oracles of Balaam (Num. 23:9, 19, 20;
24:17, see Albright, ``The Oracles of Balaam,'' \emph{passim}). The dropping
of the suffix removes the grammatical difficulties inherent in the usual
rendering of the colon; cf. note 87.

68. The \emph{mem} at the beginning of vs. 20 does not belong there (so the
LXX, Vulgate, Peshitta, and all scholars), but rather at the end of vs. 19 with
\emph{'āqēb}. It is probably to be construed as an adverbial formation (as
commonly in Ugaritic and Hebrew), ``from the rear, at the rear''; cp. Vulgate,
\emph{retrorsum}. On the other hand, the \emph{mem} may be enclitic; note that
the Peshitta omits it entirely, \emph{nqb'}.

69. See note 68.

70. MT presents a conflated text, as indicated by the disagreement in gender.
One variant is preserved in the Samaritan text, שמן לחמו (so also the LXX and
Vulgate). The other is preserved in the Peshitta and Targum; the Hebrew text
underlying these apparently read שמנה ארצו; cf. Ugar. 'Anat II:39; Gen. 27:28,
etc. The word \emph{leḥem} is used in its basic sense of ``food.''

71. Omit \emph{wəhû'}, which has been inserted in the text as a result of
dittography (note the rendering of the Vulgate). This is suggested by metrical
considerations.

72. It is also possible to read the Qal passive, \emph{yuttan}, ``he receives.''

73. The present text is unsatisfactory in many respects. Of the suggested
emendations, the best involves the change of \emph{'ayyālā} to \emph{'ēlā},
and \emph{'imrê} to \emph{'āmîrê} (cf. Skinner, \emph{Genesis}, \emph{ad
loc.}). In addition, the following readings commend themselves:
\emph{šilləḥā} for \emph{šəluḥā}, and \emph{nātənâ} for \emph{han-nōtēn},
dropping the initial \emph{he} as a case of dittography; the final \emph{â}
would not be indicated in the older orthography. The resulting text is

\begin{quote}
\begin{tabular}{ll}
Naphtali is a terebinth which produces, & נפתלי אל שלח \\
Which puts forth beautiful foliage & נתן אמר שפר \\
\end{tabular}
\end{quote}

This reconstructed text is based generally upon the LXX: Νεφθαλι στέλεχος
ἀνειμένον ἐπιδιδούς ἐν τῷ γενήματι κάλλος. This, however, does not necessarily
reflect the Hebrew text given above. For \emph{'ēlā}, we would normally expect
δρῦς or τερέβινθος, not στέλεχος, which stands for Hebrew \emph{'ālê, gēza'},
etc. In conclusion, while the proposed text is plausible, it is not entirely
convincing. In the present state of our knowledge, perhaps nothing better can
be attained. Cf. Barnes, ``A Taunt-Song in Gen. XLIX:20, 21'' \emph{JTS} 33
(1932), pp. 354--359.

74. This verse is hopelessly corrupt, as witness the versions and commentaries.
Some suggestions may be in order, however. The verse consists of a tricolon,
the first two cola of which present the familiar poetic pattern, abc:abd. The
most recent and complete treatment of this pattern in Ugaritic and Hebrew
poetry is to be found in Albright's ``The Psalm of Habakkuk,'' soon to appear
in the T. H. Robinson Anniversary volume. MT \emph{ben} is to be read with the
case ending, \emph{metri causa}, as in the Oracles of Balaam, Num. 24:3, 15;
cf. Albright, ``The Oracles of Balaam,'' p. 216, note 54. The word
\emph{pōrāt} is unintelligible as it stands, none of the suggested
interpretations having any real merit. That the text may nevertheless be
correct is indicated by the occurrence of the name \emph{prt} in Ugaritic
(cf. Gordon, \emph{Ugaritic Handbook}, Text No. 315:4), and also the
combination \emph{bn . prtn}, Text 400:III:9.

The phrase \emph{'alê 'ayin} is to be connected with \emph{'alê šûr}, on
which see the discussion below.

In the third colon, MT בנות צעדה and Samaritan בני צעירי are both in need of
correction. Samaritan בני, ``my son,'' is perhaps correct. If the \emph{taw}
in MT is taken with the following verb, it is read as a Qal imperfect,
\emph{tiṣ'ad}, ``thou dost march,'' or it may be the emphatic imperative,
\emph{ṣa'dâ}, ``march.'' In early poetry, this verb is used in a martial
sense; cf. Jud. 5:4 = Ps. 68:8. With the preposition \emph{'alê}, this suggests
a military attack. The chief targets in an attack upon a city would be the
water supply (\emph{'ayin}), and the wall (\emph{šûr}). In spite of
considerable obscurities in vss. 23--24, the same picture of military struggle
is given.

75. This verse presents no serious difficulties in itself; the chief problem
is to relate it to the surrounding material; cf. notes 74 and 76. The single
emendation of MT is taken from the Samaritan, יריבהו, \emph{yərîbûhû},
``they strive against him,'' for \emph{wārōbbû}. This improves the meter as
well as the meaning. The \emph{ba'alê ḥiṣṣîm} presumably are archers, though
the expression is unique in the Bible.

76. Verse 24a has defied the best efforts of scholars for many years. Nothing
satisfactory has resulted. The versions only add to the confusion.

77. As recognized by most scholars, MT does not make sense, and is also overly
long. In the first two words of the colon, משם רעה, there is a conflated
reading. One or two letters have dropped out by haplography, the original
variants being \emph{miššōmēr} and \emph{mērō'ê}; written out in early
orthography the text would have appeared something like this,
משמר<מא>רע or משמר<רע>. The word אבן is impossible in this context. What is
required is a parallel to \emph{ya'ăqōb}, and ``Israel-stone,'' whatever that
might mean, will not do. The simplest emendation is to read \emph{bənê} (בני)
for \emph{'eben}; the intrusion of the \emph{aleph} might be due to the
influence of אביר in the first colon of vs. 24b. On Yahweh as the
\emph{šōmēr yiśrā'ēl}, cf. Ps. 121:4; and as the \emph{rō'ê yiśrā'ēl},
Ps. 80:2. Cp. Gunkel's treatment, \emph{Genesis}, p. 427. A somewhat different
reconstruction of the text, though along similar lines, is offered by
Morgenstern in his ``Divine Triad in Biblical Mythology,'' \emph{JBL} LXIV
(1945), p 25.

78. Grammatically, the expression may be understood as either ``from the God of
thy father,'' or ``from El, thy father.'' In the present context, the former is
perhaps preferable; in support of this reading, see the parallel expression in
Ex. 15:2, אלהי אבי. On the second rendering, cf. Morgenstern, ``Divine Triad in
Biblical Mythology,'' p. 25. There are a number of references in the OT to God
as the father of his people, or of different individuals, viz. Deut. 32:6, etc.

79. Read \emph{wə'ēl}, with three Hebrew mss., Samaritan, and the Peshitta
(\emph{w'ylšdy}); the LXX has ὁ θεός ὁ ἐμός, the Vulgate, \emph{et Omnipotens}.
See, however, Morgenstern's remarks, ``Divine Triad in Biblical Mythology,''
pp. 25--26.

79a. Cf. Gen. 27:39b; Deut. 33:13, Chapter IV, note 43.

80. The bicolon (vs. 25b) is parallel to Deut. 33:13, on which see Chapter IV,
\emph{ad loc.} In its present form, the second colon here is too long, having
been harmonized with the parallel in Deut. 33. In Deut. 33:13, the word
\emph{meged} = \emph{birkôt} has been omitted; in Gen. 49, presumably one of
the other words ought to be dropped, \emph{rōbeṣet} (thus preserving the
parallelism with the first colon, and removing the metrical difficulty). In
dealing with cases of parallel transmission, the earliest written form of the
texts would exhibit the widest divergence; over the years, in the process of
transmission, the tendency would be toward harmonizing the differences. The
common original would lie much further back when the poem circulated in oral
form.

The LXX offers a variant to this colon, καὶ εὐλογίαν γῆς ἐχούσης πάντα, which
is equivalent to Deut. 33:16a. It may be a genuine variant, indicating the
parallel between heaven and earth, and not simply an inner corruption in the
transmission of the Greek mss., due to the influence of the passage in
Deuteronomy. Note that the LXX of Gen. 49:25b does not correspond at all with
that of Deut. 33:16a in so far as the wording is concerned.

This shifting of the particular blessings would point to the relatively fluid
character of the original blessing. It probably appeared in a number of forms,
of which only the two in Gen. 49 and Deut. 33 have survived. Besides the
material common to both, certain blessings are peculiar to each recension. This
suggests that the common original was somewhat longer than either, and probably
differed from both. See our remarks in Chapter IV, note 41. For this reason, it
is difficult to accept, in principle, I. Sonne's attempt to harmonize the texts
of Deut. 33:13--15 and Gen. 49:25--26, in his ``Genesis 49:25--26,'' \emph{JBL}
LXV (1946), pp. 303--306. In particular details the two texts may be useful in
correcting each other; but there is as much reason to expect differences in the
order and content of the blessings as similarities.

81. Vs. 26a, ברכת אביך גברו על, must stand in parallel with vs. 25c,
ברכת שדים ורחם. The difficulties in the text are serious, however. The best
line of approach is offered by Gunkel, who reads ברכת פת אב אף גבר ועל,
\emph{Genesis}, \emph{ad loc.} Morgenstern's rendering improves the meter by
the omission of \emph{'ābîkā}, ``Divine Triad in Biblical Mythology,''
pp. 25--26. Sonne's reading is unacceptable. On the basis of the Samaritan and
Morgenstern's correction of Gunkel's rendering, it is possible to see in MT a
defective conflated text. One variant read, in accordance with the Samaritan
and the LXX, ברכת אב ואם (the omission of אם in MT is due to homoiarkton; the
suffix is not expected, because of the general terms in the parallel colon,
and probably came in under the influence of \emph{'ābîkā} in vs. 25a;
metrical considerations are also involved). The other variant corresponds
roughly to Morgenstern's text, ברכת גבר ועל. The second variant seems to
provide a better parallel to the first colon. The first variant may originally
have formed part of another blessing.

82. With most scholars, and following Hab. 3:6, we read here
\emph{harărê 'ad}. Compare the parallel phrases in Deut. 33:15, Num. 23:7;
Chapter IV, notes 51 and 52.

83. For \emph{ta'ăwat}, read \emph{birkôt}, as indicated by the constant
repetition of \emph{birkôt} through the blessing. תאוה is presumably a
corruption of תבואת (cf. Deut. 33:14), which has been brought into the text.
See Skinner's discussion, \emph{Genesis}, p. 532. It should be pointed out that
the LXX reads εὐλογίαις = \emph{birkôt}.

84. Read *\emph{tihyan(na)}, Qal third feminine singular jussive, with a
collective subject (\emph{birkôt}); cp. the parallel Deut. 33:16, Chapter IV,
note 55. On the use of energic \emph{nun} without the pronominal suffix, see
note 67.

85. Omit the conjunction at the beginning of the second colon.

86. On the blessing as a whole, cf. Chapter IV, note 41.

87. Note that in the first colon, there is a relative clause without the
relative particle. This is a common construction in early Hebrew poetry, and
is the basis for our reconstruction of the first colon of the blessing on Gad;
cf. note 67.

88. Omit the conjunction at the beginning of cola.

89. The blessing on Benjamin is a tricolon, 3:3:3. Cp. the blessing of
Benjamin in Deut. 33:12, and Naphtali, Deut. 33:23.
